\section{Background and Related Work}
\label{sec:background}

\subsection{Skew in parallel joins}

Skew in parallel joins is old and well understood in the abstract. Walton et
al.~\cite{walton1991taxonomy} gave a taxonomy and performance model of skew
effects; DeWitt et al.~\cite{dewitt1992skew} proposed practical handling
including range partitioning with subset replication for heavy keys --- the
direct ancestor of what practitioners now call salting. In the MapReduce era,
SkewTune~\cite{kwon2012skewtune, kwon2012skewtunedemo} repartitioned a
straggler's remaining input at runtime, and Mantri~\cite{ananthanarayanan2010mantri}
established at production scale that outlier tasks dominate job latency.

A substantial line of work proposes new partitioners. For Spark specifically:
intermediate-data partitioning for skew mitigation~\cite{tang2021intermediate}
and skew-resilient adaptive parallel processing~\cite{shen2020srspark}. In the
MapReduce lineage from which these descend: sampling-based and cost-aware
partitioning schemes~\cite{myung2016handling, gufler2011skew}. Irandoost et
al.~\cite{irandoost2019skewslr} survey the MapReduce field systematically.
Their review discusses neither Adaptive Query Execution nor salting --- fairly
enough for AQE, which shipped in Spark 3.0 after that review was written, but
the omission is indicative: the literature has been optimising the mechanism
while the question of \emph{selecting} among deployed mechanisms went unasked.

Metwally's work on scaling and load-balancing equi-joins~\cite{metwally2025amjoin}
is the closest formal treatment of the idea salting embodies. AM-Join combines
randomised key expansion with replication on the opposite side and analyses
doubly-skewed keys as the hard case. It is the right formal ancestor to cite,
and it is not a study of the technique practitioners actually run.

\subsection{Adaptive Query Execution}

AQE re-optimises between stages using runtime statistics: coalescing
post-shuffle partitions, converting sort-merge joins to broadcast joins when
the observed side turns out small, and splitting skewed
partitions~\cite{databricks_aqe, spark_aqe_docs}. Xue et
al.~\cite{xue2024adaptive} describe the production descendant of this machinery
at lakehouse scale and state its skew handling explicitly: the rule discovers
skew on the join keys of a shuffled hash join, ``which manifests as a few
partitions containing significantly more data than others,'' and eliminates it
by ``logically splitting those large, consumer-side partitions into smaller,
more balanced partitions.''

That paper is peer-reviewed, and it is the authoritative account of the
mechanism. It is also a systems-description paper written by the vendor and
evaluated on the vendor's engine. What it does not provide --- and what does not
exist elsewhere --- is an independent, controlled evaluation of the open-source
rule against the alternative practitioners actually reach for.

Adjacent work confirms the shape of the gap rather than filling it. Learned and
adaptive Spark optimisers --- AQORA~\cite{he2025aqora},
LOCAT~\cite{xin2022locat}, fine-grained parameter tuning~\cite{lyu2024sparkopt},
zero-execution configuration tuning~\cite{suri2025zeroexec} --- address join
ordering, broadcast selection and configuration knobs. AQORA operates
\emph{inside} the AQE framework and does not address skew. Zhang et
al.~\cite{zhang2023adaptive} evaluate adaptive query processing rigorously, but
on single-node relational engines rather than Spark.

\subsection{The specific absence}

We searched dblp, the arXiv full-metadata API, Crossref, and domain-scoped
searches over the ACM Digital Library, IEEE Xplore, the VLDB proceedings,
Springer, ScienceDirect and MDPI. Two negative results are load-bearing for
this paper, and we state them with their limits.

\textbf{No independent evaluation of Spark's skew-join rule.} We found no
peer-reviewed paper whose contribution is an evaluation of
\texttt{spark.sql.adaptive.skewJoin}. The nearest comparative study, Phan et
al.~\cite{phan2022skewjoin}, compares hash-based, range-based and
multi-dimensional range partitioning on MapReduce and Spark --- and does not
consider AQE. The mechanism is otherwise documented only in vendor
material~\cite{databricks_aqe, spark_aqe_docs} and described from the
provider's perspective~\cite{xue2024adaptive}.

\textbf{No peer-reviewed treatment of salting.} A full-metadata arXiv search
for \texttt{abs:salting AND abs:join} returns five papers, all in
condensed-matter physics and astronomy. Searches for ``key salting'' and
``salted join'' across the major publisher databases return only unrelated
hash-join work. Every source describing the technique as practitioners apply it
is a blog post, a tutorial, or vendor documentation. Its formal analogue is
randomised key expansion with replication~\cite{metwally2025amjoin,
dewitt1992skew}, never under that name and never as the object of study.

We note the limits of this claim. Semantic Scholar and OpenAlex were
unavailable during our search, so recall rests on dblp, arXiv, Crossref and
domain-scoped web search, and dblp's coverage of some journals is imperfect ---
Phan et al. is itself not indexed there. We therefore claim these absences as
strong evidence rather than as proof.

\subsection{Cost models and benchmarking practice}

Cost modelling for big-data query processing is well
developed~\cite{siddiqui2020costmodels, zhou2010scopeopt, leis2021costoptimal},
including for shuffle specifically~\cite{shen2020magnet}, but we are not aware
of a cost model for the salting decision. Methodologically we follow
LST-Bench~\cite{camacho2024lstbench} in building the benchmark from composable,
declared workload patterns, and Nurdin et al.~\cite{nurdin2026lakehouseperf} in
treating runtime variance as a first-class threat: they measure roughly
twofold run-to-run variation for identical lakehouse queries on shared cloud
compute, which is precisely why our timed measurements are taken on dedicated
local compute and our transferable claims rest on deterministic counters.
